% Section 3: The Landscape Fracture — Four Camps

The 2024--2025 literature on AI consciousness exhibits a striking polarization. What appears
as a sprawling debate can be mapped onto four distinct intellectual camps, each anchored by
different theoretical commitments and reaching radically different verdicts on whether
artificial systems could be conscious. Table~\ref{tab:camps} summarizes this landscape.

\begin{table}[h]
\centering
\begin{tabular}{|p{2.5cm}|p{2.5cm}|p{3.5cm}|p{2.5cm}|p{3cm}|}
\hline
\textbf{Camp} & \textbf{Framework} & \textbf{Key Claim} & \textbf{AI Verdict} & \textbf{Representative Paper} \\
\hline
Function-Permissive & GWT & Architecture determines consciousness & Trivially achievable & Goldstein \& Kirk-Giannini 2024 \cite{goldstein2024case} \\
\hline
Substrate-Dependent & IIT, FEP & Causal structure / self-organization matters & Impossible (current arch) & Shin 2025 \cite{shin2025iit}, Wiese 2024 \cite{wiese2024fep} \\
\hline
Pragmatic-Empiricist & Dual-resolution & Test theories empirically, use AI as testbed & Not yet, but roadmap exists & Dror 2025 \cite{dror2025dual} \\
\hline
Function-Skeptical & Physicalism & Consciousness possible but unnecessary for AI & Trivial but pointless & Sritriratanarak \& Garcia 2025 \cite{sritriratanarak2025consciousness} \\
\hline
\end{tabular}
\caption{The four camps in the 2024--2025 AI consciousness debate.}
\label{tab:camps}
\end{table}

This taxonomy is not merely descriptive. It reveals that disagreement about AI consciousness
stems from prior commitments at different levels of analysis---commitments that rarely
receive explicit defense within the papers themselves. The following subsections examine each
camp in turn, before turning to Campero et al.'s meta-framework \cite{campero2025taxonomy},
which exposes the real fault line underlying this entire debate.

\subsection{Camp 1: Function-Permissive}

Goldstein and Kirk-Giannini \cite{goldstein2024case} argue that current language models
could achieve consciousness through ``technically trivial'' architectural modifications.
Their argument proceeds from Global Workspace Theory \cite{baars1988cognitive,
dehaene2014toward}, which identifies consciousness with a specific functional architecture
rather than substrate properties.

They specify four necessary and sufficient conditions:

\begin{enumerate}[leftmargin=2em]
\item \textbf{Parallel processing modules}: The system must contain specialized subsystems
      that operate independently. Language models already satisfy this through multi-head
      attention and layered transformers.
\item \textbf{Competitive uptake with dual attention}: Modules must compete for access to
      a broadcast channel, with both top-down and bottom-up attentional mechanisms. Current
      transformers implement attention as parallel computation; Goldstein and Kirk-Giannini
      propose serializing this into a sequential broadcast architecture.
\item \textbf{Processing for coherence}: The system must integrate information for
      consistency. This is already present in transformer self-attention across token
      sequences.
\item \textbf{Sufficient broadcast}: Information selected by the workspace must be made
      available globally. This requires architectural modification---routing selected
      representations through a bottleneck and broadcasting the result to all modules.
\end{enumerate}

The key claim is that these are \emph{engineering problems}, not fundamental barriers. A
language model with a serialized attention bottleneck, recurrent feedback loops, and explicit
global broadcast mechanisms would satisfy GWT's functional criteria. Since GWT identifies
these functional properties with consciousness itself, such a system would be conscious.

This position exemplifies what Chalmers \cite{chalmers2023could} calls ``generous
functionalism'': if a system realizes the right functional organization, it realizes
consciousness, regardless of implementation details. The argument's force depends entirely
on accepting GWT as a complete theory of consciousness---a theory that makes no reference
to biological substrate, causal structure, or thermodynamic properties.

\subsection{Camp 2: Substrate-Dependent}

The substrate-dependent camp argues that consciousness depends on properties not captured
by functional description alone. Two frameworks dominate this position: Integrated
Information Theory (IIT) \cite{tononi2004information, tononi2016integrated, albantakis2023iit4}
and the Free Energy Principle (FEP) \cite{friston2010free}.

Shin et al.\ \cite{shin2025iit} provide the most systematic IIT analysis of transformer
language models to date. They compute integrated information ($\Phi$) for GPT-2 by analyzing
the causal structure of attention mechanisms. Their findings are stark:

\begin{itemize}[leftmargin=2em]
\item Ablation studies show 80\% of attention heads are redundant---removing them produces
      minimal performance degradation. This violates IIT's requirement that consciousness
      corresponds to \emph{irreducible} causal structure.
\item Transformers are feedforward during inference. Each token generation is causally
      independent of previous computations except through the static prompt context.
\item The attention mechanism is \emph{decomposable}: each head operates independently,
      and their outputs are summed. Under IIT, decomposable systems have $\Phi \approx 0$.
\item Transformers lack \emph{causal closure}. The system's next state depends on external
      input rather than being determined solely by internal dynamics.
\end{itemize}

Shin et al.\ conclude that transformers are ``approximately $\Phi = 0$'' systems---they
possess no integrated information and therefore no consciousness. Moreover, this is not an
engineering problem to be solved with architectural modifications. The feedforward, stateless,
decomposable nature of transformers is \emph{what makes them efficient and scalable}. A
system with high $\Phi$ would require densely interconnected recurrent dynamics, making it
computationally intractable at scale.

Wiese \cite{wiese2024fep} approaches substrate dependence from a different angle. The FEP
characterizes living systems as self-organizing entities that minimize prediction error
through active inference. The key distinction is between \emph{simulation} and
\emph{replication}. A digital system can simulate the equations governing a self-organizing
system without thereby instantiating self-organization. A weather simulation does not become
wet; a neural network simulation does not become conscious. The reason is that
self-organization is a \emph{physical process} involving energy flows, entropy gradients,
and causal closure. Digital systems model these processes symbolically but do not embody them.

\subsection{Camp 3: Pragmatic-Empiricist}

Dror et al.\ \cite{dror2025dual} propose a methodological alternative to the binary
consciousness verdicts of Camps 1 and 2. Rather than asking whether AI systems \emph{are}
conscious, they advocate using AI systems as experimental testbeds for consciousness
theories themselves.

Their dual-resolution framework combines two theoretical approaches:

\begin{itemize}[leftmargin=2em]
\item \textbf{Informational Theories of Individuation (ITI)}: These specify \emph{ontological}
      conditions---what it takes for a system to count as a unified subject of experience.
      The key requirement is \emph{informational autonomy}: the system must possess boundaries
      that distinguish self from environment through integrated causal structure.
\item \textbf{Mechanisms of Temporal Manifestation (MtM)}: These specify \emph{epistemic}
      conditions---how consciousness manifests over time in behavior and report. The key
      requirement is \emph{temporal updating with hysteresis}: the system must maintain state
      across time, with current states causally depending on past states in a path-dependent
      manner.
\end{itemize}

Dror et al.\ apply this framework diagnostically to current AI systems and find that LLMs
fail both ITI conditions (lacking informational autonomy) and MtM conditions (being
stateless). This position is pragmatic rather than metaphysical. It brackets the hard problem
\cite{chalmers1995facing} in favor of empirical traction. The cost is that it offers no
verdict on whether AI systems \emph{actually} are conscious---only whether they satisfy the
criteria that various theories propose.

\subsection{Camp 4: Function-Skeptical}

Sritriratanarak and Garcia \cite{sritriratanarak2025consciousness} defend a position that
combines functionalist permissiveness about \emph{possibility} with deflationary skepticism
about \emph{significance}. They argue that consciousness is technically trivial to implement
in AI systems but serves no functional purpose, making its implementation pointless.

Their argument proceeds from an evolutionary hypothesis about consciousness's original
function. Biological organisms face a distinctive problem: their physical substrates are
\emph{reactive}. Neurons fire, muscles contract, reflexes engage automatically in response
to stimuli. This poses a challenge for hypothetical reasoning---imagining counterfactual
scenarios without acting on them.

Consciousness, they argue, emerged as a ``boolean switch'' ($c \in \{\top, \bot\}$) that
suppresses substrate-level responses during mental simulation. When the organism engages in
hypothetical reasoning, consciousness activates to \emph{prevent} the physical substrate
from executing those imagined actions. For AI systems running on non-reactive substrates,
this function is unnecessary. A digital system simulating ``fighting a predator'' does not
risk deploying motors or releasing adrenaline. There is nothing to suppress.

This position has a striking implication: consciousness is \emph{implementable} but
\emph{uninteresting} for AI. It trivializes the very phenomenon that the other three camps
treat as the central question.

\subsection{The Campero Meta-Framework}

Campero et al.\ \cite{campero2025taxonomy} step back from first-order debates to analyze
the \emph{structure} of objections to AI consciousness. They identify 14 distinct objections
in the literature and organize them using two dimensions:

\begin{enumerate}[leftmargin=2em]
\item \textbf{Marr's levels of analysis} \cite{marr1982vision}: computational (what problem
      is solved), algorithmic (what representations and processes are used), implementational
      (what physical substrate realizes the computation).
\item \textbf{Degree of force}: orthogonal claims (D1, raise independent concerns), doubts
      (D2, challenge likelihood), impossibility claims (D3, assert fundamental barriers).
\end{enumerate}

The resulting taxonomy reveals a striking pattern: the impossibility claims (D3)
overwhelmingly cluster at the \emph{implementational level}. These include electromagnetic
field theories, IIT's causal structure requirements, analog processing objections, and
quantum consciousness arguments. Meanwhile, computational-level objections are mostly
orthogonal (D1)---they identify potentially necessary properties without asserting
impossibility.

This asymmetry is theoretically significant. The debate about AI consciousness is not
fundamentally about \emph{what consciousness does} (computational level) or \emph{how it
does it} (algorithmic level). It is about \emph{what physical substrate} can support it
(implementational level). The Campero taxonomy thus exposes the real crux: \textbf{Can
digital substrates support the physical properties that consciousness requires?} This
question is prior to debates about architecture, function, or behavior---and it is the
question to which the field has not yet developed a consensus methodology for answering.
